% §3 Methodology, v3 (rewritten 2026-06-08 for the 12-city / AOAS reframe).
% Replaces the old 3-city methods built around GRL adversarial deconfounding,
% binarized DR, a randomization intervention, and SCM0/SCM1 cross-market
% falsification. Describes the methods the v3 results actually use: partialling-out
% DML, the Shen and Baur text channels, LEACE+IGBP erasure as a diagnostic, the
% Gelbach fixed-treatment decomposition (headline test), the LLM counterfactual,
% Cinelli-Hazlett sensitivity, and random-effects pooling. Style: flowing prose,
% no em dashes, no enumerations, no punchy fragments.

\section{Methodology}
\label{sec:methods}

\subsection{Setup and the conflation it invites}

Write $\mathcal{D} = \{(x_i, t_i, \ell_i, y_i)\}_{i=1}^{n}$ for a market of $n$
listings, where $x_i \in \mathbb{R}^{d}$ collects the structured property
attributes that hedonic models have always used, the bedroom and bathroom counts,
the living area, the lot size, and the year built; $\ell_i \in \mathbb{R}^{2}$ is
the latitude-longitude coordinate of the parcel; $y_i$ is the log sale price; and
$t_i$ is a representation of the listing's free-text description, either a scalar
summary or a high-dimensional embedding depending on the channel. A text-augmented
valuation model learns some $\hat{y}_i = f(x_i, t_i, \ell_i)$, and the empirical
regularity that motivates this literature is that adding $t_i$ to the feature set
raises predictive accuracy beyond what $x_i$ and $\ell_i$ deliver on their own.

The difficulty is that an accuracy gain cannot say what kind of information the
text carried. An agent who writes ``sun-drenched and steps from the park'' is
describing a property, but the same phrase is also a near-deterministic signal of
which block the property sits on, because language and location are produced
together by an author who observes both. A predictive model that prices genuine
semantic content and a predictive model that prices laundered geography make
identical out-of-sample predictions, so the gain on its own leaves an analyst
unable to tell which of the two has been bought. Our object of interest is the
part of the text-price relationship that is not the spatial information the text
carries, and the methods that follow are built to isolate it and then to ask, for
each metropolitan market in turn, how much of the apparent text signal that
isolation leaves standing.

\subsection{A structural account of the confounding}

The data-generating process we have in mind makes the confounding explicit rather
than assuming it away. Location is exogenous and drives everything downstream: the
structured attributes $x$, a vector $c$ of contextual neighborhood features such
as demographics, amenities, and recent crime, the text $t$, and the price $y$. The
text in turn is written with knowledge of location, structure, and context, so
$t = f_t(\ell, x, c, U_t)$, and the question the paper exists to answer is whether
$t$ enters the structural equation for $y$ at all once $\ell$, $x$, and $c$ are
accounted for, or whether the entire observed association between text and price
runs through the common dependence of both on location.

\begin{figure}[!htbp]
\centering
\begin{tikzpicture}[
    node distance=1.5cm and 2cm,
    every node/.style={draw, circle, minimum size=0.8cm},
    arrow/.style={->, >=stealth, thick}
]
    \node (L) {$\ell$};
    \node[right=of L] (X) {$x$};
    \node[right=of X] (C) {$c$};
    \node[below=of X] (T) {$t$};
    \node[right=of C] (Y) {$y$};

    \draw[arrow] (L) -- (X);
    \draw[arrow] (L) to[bend left=20] (C);
    \draw[arrow] (L) to[bend left=30] (T);
    \draw[arrow] (X) -- (T);
    \draw[arrow] (C) -- (T);
    \draw[arrow] (L) to[bend left=15] (Y);
    \draw[arrow] (X) -- (Y);
    \draw[arrow] (C) -- (Y);

    \draw[arrow, red, very thick, opacity=0.4] (L) to[bend left=30] (T);
    \draw[arrow, red, very thick, opacity=0.4] (T) -- (Y);
\end{tikzpicture}
\caption{The confounding structure the paper adjudicates. Location $\ell$ is
exogenous and feeds the property attributes $x$, the contextual features $c$, the
listing text $t$, and the price $y$. The shaded path $\ell \rightarrow t$ together
with $\ell \rightarrow y$ is the spatial-confounding concern: text and price are
associated in part because both depend on where the property is, so the relevant
causal quantity is the direct $t \rightarrow y$ effect that survives conditioning
on $\ell$, $x$, and $c$. Whether that direct effect is zero or positive is an
empirical question the markets answer market by market.}
\label{fig:dag}
\end{figure}
\FloatBarrier

Two readings of this diagram bracket the possibilities. Under one, the text effect
on price is zero net of location and the structured covariates, every observed
text-price correlation is geography wearing a semantic costume, and the apparent
predictive value of listing language is spatial confounding mediated through words.
Under the other, listing text carries property-specific content that the
structured covariates miss, so $t$ does enter the equation for $y$ and the
text-price relationship is at least partly a genuine semantic effect. The two
readings disagree about a quantity that is identified once the backdoor set
$\{\ell, x, c\}$ is conditioned on, since that set blocks every path from $t$ to
$y$ that does not pass through the hypothesized direct edge, and the rest of this
section develops the estimator, the encodings, and the deconfounding tests through
which the twelve markets adjudicate between them.

\subsection{The partially-linear estimand and its estimation}

We summarize the listing text by a scalar treatment $T$, defer to
Section~\ref{sec:methods-channels} how that scalar is constructed, and pose the
partially-linear model
\begin{equation}
Y = \theta\,T + g(\ell, x, c) + \varepsilon, \qquad
T = m(\ell, x, c) + \nu,
\label{eq:plr}
\end{equation}
in which $\theta$ is the per-standard-deviation effect of the text direction on
log price and the nuisance functions $g$ and $m$ are left unrestricted. The
Robinson partialling-out estimator \citep{robinson1988root} recovers $\theta$ from
the residual-on-residual regression of $\tilde{Y} = Y - \hat{g}$ on
$\tilde{T} = T - \hat{m}$, and we fit $\hat{g}$ and $\hat{m}$ by cross-fitted
gradient boosting in the double machine learning framework of
\citet{chernozhukov2018double}, splitting each market into $K = 5$ folds and
predicting each fold's residuals from models trained on the complement so that the
nuisance estimates and the score are never formed on the same observations. The
resulting estimator is Neyman-orthogonal and root-$n$ consistent whenever the
nuisance fits converge faster than $n^{-1/4}$, and its influence-function variance
gives the standard error we report by default,
\begin{equation}
\hat{\theta} = \frac{\sum_i \tilde{T}_i \tilde{Y}_i}{\sum_i \tilde{T}_i^{2}},
\qquad
\hat{\mathrm{SE}}_\theta
= \sqrt{\frac{1}{n^{2}} \sum_i
\frac{(\tilde{Y}_i - \hat{\theta}\,\tilde{T}_i)^{2}\,\tilde{T}_i^{2}}
{\big(\tfrac{1}{n}\sum_j \tilde{T}_j^{2}\big)^{2}}}.
\label{eq:dml}
\end{equation}

One feature of the embedding channel complicates the inference and we flag it here
because it shapes how the per-market intervals should be read. When the treatment
is the leading principal component of the embedding, that direction is itself
estimated on the same sample that estimates the price effect, so $T$ is a generated
regressor in the sense of \citet{pagan1984econometric} and the orthogonal-score
variance in Equation~\eqref{eq:dml} does not absorb the extra sampling variability
the estimation of the direction introduces. The interval is correspondingly
anticonservative at the few thousand listings a single market affords, mildly so
near the null and more seriously when the true effect is large, a pattern we
document in a calibration study and repair with an out-of-fold construction of the
treatment direction in Appendix~\ref{app:simulation}. The point estimates are
approximately unbiased throughout; it is the nominal coverage of the analytic
interval, not the location of the estimate, that the calibration study qualifies.

\subsection{Two encodings of the text treatment}
\label{sec:methods-channels}

Because no single reduction of a paragraph to a number is obviously the right one,
we read the listing text through two encodings with different inductive biases and
treat their agreement as a robustness check rather than as independent
corroboration. The first follows the hedonic specification of
\citet{shen2021information}, who summarize a description by a Doc2Vec uniqueness
score that measures how far a listing's language departs from the corpus norm, and
we enter that scalar into Equation~\eqref{eq:plr} as the treatment. The second
follows \citet{baur2023automated}, who embed each description with a pretrained
sentence transformer and read the leading principal component of the embedding as
the text signal; we standardize that component and use it as the treatment in a
separate partialling-out fit. The sign of a principal component is not identified,
since a component and its negation span the same axis, so wherever the Baur channel
appears we orient the common axis to make the pooled estimate non-negative and
disclose the orientation as a convention rather than a finding, a convention that
cannot manufacture agreement because the cross-market correlation between channels
is invariant in magnitude under a sign flip of either one.

\subsection{What the embedding encodes: erasure as a diagnostic}

The sentence-transformer embedding that underlies the Baur channel was pretrained
on generic text and has no reason to keep semantic content about a property
separate from the geographic information that correlates with its neighborhood, so
before asking whether the price effect is confounded we ask the prior question of
how much location the representation carries at all. We answer it by erasing the
coordinate concept from the embedding and measuring what an erasure leaves
recoverable. The closed-form concept erasure of \citet{belrose2023leace} removes
the linear subspace that the mean-conditional-covariance estimator associates with
the latitude-longitude pair and comes with an exact-guardedness guarantee that no
linear predictor can recover the coordinate from the residual better than a
constant baseline, and the iterative nonlinear procedure of
\citet{iskander2023shielded} extends the attack to the directions a linear probe
cannot see by alternating between fitting a nonlinear probe of the coordinate and
projecting out where that probe concentrates its predictive mass. The erasure is a
diagnostic of what the representation contains and not a modification of any
per-market estimate, since the Baur effects are computed on the unprojected
embedding in keeping with the published protocol; its role is to bound how much
location a linear or nonlinear function of the text could carry, which the
decomposition below then connects to how much location actually moves the price.

\subsection{Holding the text fixed: a covariate decomposition}

Showing that the embedding encodes geography is not the same as showing that the
price effect is confounded by it, because a representation can carry a confounder
along directions that have nothing to do with the one on which the treatment acts.
The test that speaks to confounding directly holds the text treatment fixed and
moves geography across the conditioning set, in the manner of the order-invariant
covariate decomposition of \citet{gelbach2016covariates}. We estimate the text
effect twice on the same fixed treatment, once with the location coordinates held
out of the controls and once with a flexible spatial basis, the coordinates
together with their squares and their interaction, folded in, and read the
difference between the two estimates as the share of the text effect that explicit
geography accounts for. A small difference means the text-price effect is not the
spatial information re-expressed, and the decomposition is informative rather than
mechanical precisely when the text direction is shown by the erasure diagnostic to
carry geography, since a direction that encodes location yet whose price effect is
unmoved by conditioning on location is the signature of a genuine semantic signal.
This contrast is the paper's headline deconfounding test, and
Section~\ref{sec:decomposition} reports it for every market.

\subsection{An intervention check and a sensitivity bound}

Two further instruments probe the same question from angles with different failure
modes. A counterfactual rewriting intervention uses an open-weight language model
to edit each description under a minimal-pair protocol, producing one rewrite that
strips stylistic content while holding the model's inferred submarket fixed and a
second that also moves the targeted submarket, and the price differentials these
edits imply isolate a within-listing style contrast and a location-inference
contrast whose biases are not shared with the cross-listing geometry of a fixed
encoder. Because none of these designs identifies a structural effect by
assumption, we close with the sensitivity calculus of \citet{cinelli2020making},
reporting for each market the robustness value, the strength an unobserved
confounder would need on both the treatment and the outcome to overturn the
estimate, benchmarked against the strength of the spatial covariates we do observe.
The twelve per-market estimates are finally combined through a random-effects
meta-analysis with the Paule-Mandel variance estimator and Hartung-Knapp-Sidik-Jonkman
small-sample intervals \citep{paule1982consensus, inthout2014hartung}, and a
moderator scan asks, with multiplicity and leverage diagnostics, whether any
metropolitan covariate explains the heterogeneity the markets display.
